\begin{corollary}\label{cor:kahler-smooth-dual-attainment}
Let $\omega$ be K\"ahler on a connected compact complex manifold of dimension $n\geq2$. For a nonzero nonnegative $f\in L^1(\mathrm{d}V)$, the functional $D_{f^{1/n}}$ has a smooth strictly positive maximizer in $\mathscr G_\omega$ if and only if $f$ agrees almost everywhere with a smooth strictly positive function.
\end{corollary}
\begin{proof}
The K\"ahler formula in \cref{prop:individual-normalizer-lp} makes $b_\omega(f)>0$. If a smooth strictly positive maximizer exists, the proof of \cref{cor:smooth-dual-exact-equation} gives
\[
 \omega+\ddc u=a\gamma,
 \qquad a\geq0,
\]
with $a,u$ smooth and $a\not\equiv0$. Since $\mathrm{d}\omega=0$, the first-order equation in \cref{lem:smooth-conformal-current} becomes $\mathrm{d}a=-aC$ for a smooth one-form $C$. Along any smooth path $\sigma$, this is the scalar equation
\[
 (a\circ\sigma)'=-(a\circ\sigma)C(\dot\sigma).
\]
Its solution is the initial value multiplied by a positive exponential. Connectedness therefore gives $a>0$ everywhere, and $f=b_\omega(f)^{-1}a^n\gamma^n/\mathrm{d}V$ is smooth and strictly positive. Conversely, the smooth equation and \cref{thm:gaud-volume-dual-formula} supply the maximizing Gauduchon form.
\end{proof}

The conclusion about strict positivity changes for a nonclosed background, even on a K\"ahler manifold.

\begin{example}\label{ex:nonclosed-dual-zero-density}
On every connected compact complex manifold of dimension $n\geq2$, there is a nonclosed Hermitian form $\omega$ and a nonzero smooth nonnegative density $f$ vanishing on an open set such that $b_\omega(f)=1$ and $D_{f^{1/n}}$ has a smooth strictly positive maximizer.
\end{example}
\begin{proof}
Fix a Gauduchon metric $\gamma$ and a point $x$. Choose a smooth real function $\psi$ such that $\ddc\psi\leq-\gamma$ on a coordinate neighborhood $U$ of $x$. One obtains it by multiplying $-M|z|^2$ by a cutoff that equals one near $x$, with $M$ large enough, and then shrinking $U$.

Choose a smooth function $\chi\geq0$ which vanishes on a smaller neighborhood of $x$, equals one outside $U$, and has values at most one. Since $X$ is compact, for sufficiently large $L>0$ the form
\[
 a=L\chi,\qquad \omega=a\gamma-\ddc\psi
\]
is strictly positive: on $U$ the second term is at least $\gamma$, and outside $U$ the term $L\gamma$ dominates $\ddc\psi$. Set $f=a^n\gamma^n/\mathrm{d}V$. The smooth potential $\psi$ satisfies $\omega+\ddc\psi=a\gamma\geq0$ and $\mathrm{MA}_\omega(\psi)=f\,\mathrm{d}V$, so $b_\omega(f)\geq1$.

Stokes and $\ddc\gamma^{n-1}=0$ give
\[
 \int_X\omega\wedge\gamma^{n-1}
 =\int_Xa\gamma^n
 =\int_X(f\,\mathrm{d}V/\gamma^n)^{1/n}\gamma^n.
\]
The Gauduchon formula \cref{thm:integrable-gauduchon-formula} gives the opposite inequality $b_\omega(f)\leq1$. The displayed test thus attains equality, and its normalized current is a smooth strictly positive maximizer. If $\omega$ were closed, the preceding corollary would force $f$ to be strictly positive, contrary to its vanishing on an open set. Thus $\omega$ is nonclosed.
\end{proof}

The restriction $n\geq2$ in \cref{cor:smooth-dual-exact-equation} is also necessary. On a compact connected Riemann surface, a normalized Gauduchon current has degree zero and is the constant $(\int_X\omega)^{-1}$: the equation $\ddc T=0$ makes it smooth and harmonic. This single smooth strictly positive current maximizes for every nonzero nonnegative $L^1$ density, which need not have a smooth representative.
